\begin{trapbox}
\textbf{误区1：用 $S<1$ 保证存在。}
$0<S<1$ 时候选线根本不相交；$S>1$ 时反而存在同支弦。
在 $y_0\ne0$ 下必须使用 $S<0$ 或 $S>1$。
\par\medskip\noindent
\textbf{误区2：$y_0=0$ 就没有弦。}
这时要另行讨论：若 $|x_0|>a$，竖直线 $x=x_0$ 有一条以 $(x_0,0)$ 为中点的弦。
若 $M=O$，过原点的直线 $y=kx$ 在 $|k|<b/a$ 时给出中心弦。
其余 $y_0=0$ 的点不能作为非退化弦的中点。
\par\medskip\noindent
\textbf{误区3：求出轨迹方程便结束。}
必须排除不满足实交点条件的部分；$S=0$ 对应渐近线方向，$S=1$ 对应重合端点。
\end{trapbox}
